\documentclass[10pt,a4paper]{amsart}
\usepackage{amsmath,amssymb,amsthm,graphicx}
% Retained minimal preamble for the Stacks Project chapter more-algebra.tex (More on Algebra).
% The chapter file pinned-original-source.tex carries only \input{preamble} on its first line, so no
% preamble text lives inside the chapter itself.  The Stacks Project uses one single shared preamble for
% every chapter; the snapshot stages it identically as sources/stacks-chow/preamble.tex and
% sources/stacks-etale-cohomology/preamble.tex.  The block below is copied mechanically, line for line,
% from sources/stacks-chow/preamble.tex lines 14-18 (the Xy-pic setup) and 154-217 (the native theorem
% declarations together with the mathematical macros).  Nothing was added, removed or redefined; the
% wrapper class and the amsmath/amssymb/amsthm/graphicx base packages are supplied by the assembler.
% Chapter: more-algebra.tex -- More on Algebra.
% For commutative diagrams we use Xy-pic
\usepackage[all]{xy}

% We use 2cell for 2-commutative diagrams.
\xyoption{2cell}

% Theorem environments.
%
\theoremstyle{plain}
\newtheorem{theorem}[subsection]{Theorem}
\newtheorem{proposition}[subsection]{Proposition}
\newtheorem{lemma}[subsection]{Lemma}

\theoremstyle{definition}
\newtheorem{definition}[subsection]{Definition}
\newtheorem{example}[subsection]{Example}
\newtheorem{exercise}[subsection]{Exercise}
\newtheorem{situation}[subsection]{Situation}

\theoremstyle{remark}
\newtheorem{remark}[subsection]{Remark}
\newtheorem{remarks}[subsection]{Remarks}

\numberwithin{equation}{subsection}

% Macros
%
\def\lim{\mathop{\mathrm{lim}}\nolimits}
\def\colim{\mathop{\mathrm{colim}}\nolimits}
\def\Spec{\mathop{\mathrm{Spec}}}
\def\Hom{\mathop{\mathrm{Hom}}\nolimits}
\def\Ext{\mathop{\mathrm{Ext}}\nolimits}
\def\SheafHom{\mathop{\mathcal{H}\!\mathit{om}}\nolimits}
\def\SheafExt{\mathop{\mathcal{E}\!\mathit{xt}}\nolimits}
\def\Sch{\mathit{Sch}}
\def\Mor{\mathop{\mathrm{Mor}}\nolimits}
\def\Ob{\mathop{\mathrm{Ob}}\nolimits}
\def\Sh{\mathop{\mathit{Sh}}\nolimits}
\def\NL{\mathop{N\!L}\nolimits}
\def\CH{\mathop{\mathrm{CH}}\nolimits}
\def\proetale{{pro\text{-}\acute{e}tale}}
\def\etale{{\acute{e}tale}}
\def\QCoh{\mathit{QCoh}}
\def\Ker{\mathop{\mathrm{Ker}}}
\def\Im{\mathop{\mathrm{Im}}}
\def\Coker{\mathop{\mathrm{Coker}}}
\def\Coim{\mathop{\mathrm{Coim}}}

% Boxtimes
%
\DeclareMathSymbol{\boxtimes}{\mathbin}{AMSa}{"02}

%
% Macros for moduli stacks/spaces
%
\def\QCohstack{\mathcal{QC}\!\mathit{oh}}
\def\Cohstack{\mathcal{C}\!\mathit{oh}}
\def\Spacesstack{\mathcal{S}\!\mathit{paces}}
\def\Quotfunctor{\mathrm{Quot}}
\def\Hilbfunctor{\mathrm{Hilb}}
\def\Curvesstack{\mathcal{C}\!\mathit{urves}}
\def\Polarizedstack{\mathcal{P}\!\mathit{olarized}}
\def\Complexesstack{\mathcal{C}\!\mathit{omplexes}}
% \Pic is the operator that assigns to X its picard group, usage \Pic(X)
% \Picardstack_{X/B} denotes the Picard stack of X over B
% \Picardfunctor_{X/B} denotes the Picard functor of X over B
\def\Pic{\mathop{\mathrm{Pic}}\nolimits}
\def\Picardstack{\mathcal{P}\!\mathit{ic}}
\def\Picardfunctor{\mathrm{Pic}}
\def\Deformationcategory{\mathcal{D}\!\mathit{ef}}

% --- appended 2026-10-12 (Chief): the project own definitions of the reference/slogan/history
% display environments, copied mechanically from the shared Stacks preamble lines 8-12
% (\usepackage{verbatim} plus the three \newenvironment lines, which the chapter preamble keeps
% verbatim from the project preamble). Nothing added, removed or redefined: these environments
% carry no mathematical content and the project itself typesets them as comments. Without these
% three lines a counted statement whose own text uses them fails with Environment slogan undefined.
\usepackage{verbatim}
\newenvironment{reference}{\comment}{\endcomment}
\newenvironment{slogan}{\comment}{\endcomment}
\newenvironment{history}{\comment}{\endcomment}
\usepackage{cleveref}
\crefname{originaldefinition}{original definition}{original definitions}
\Crefname{originaldefinition}{Original definition}{Original definitions}
\makeatletter
\newlabel{definition-pseudo-coherent}{{15874}{source}{}{sourcepointer-1}{}}
\newlabel{definition-pseudo-coherent@cref}{{[originaldefinition][0][]15874}{[1][0][]source}{}{}{}}
\makeatother
\crefname{originallemma}{original lemma}{original lemmas}
\Crefname{originallemma}{Original lemma}{Original lemmas}
\makeatletter
\newlabel{derived-remark-check-on-generator}{{0ATH}{source}{}{sourcepointer-2}{}}
\newlabel{derived-remark-check-on-generator@cref}{{[originallemma][0][]0ATH}{[1][0][]source}{}{}{}}
\makeatother
\crefname{originallemma}{original lemma}{original lemmas}
\Crefname{originallemma}{Original lemma}{Original lemmas}
\makeatletter
\newlabel{lemma-RHom-out-of-projective}{{18708}{source}{}{sourcepointer-3}{}}
\newlabel{lemma-RHom-out-of-projective@cref}{{[originallemma][0][]18708}{[1][0][]source}{}{}{}}
\makeatother
\crefname{originallemma}{original lemma}{original lemmas}
\Crefname{originallemma}{Original lemma}{Original lemmas}
\makeatletter
\newlabel{lemma-bounded-below-tor-amplitude}{{16930}{source}{}{sourcepointer-4}{}}
\newlabel{lemma-bounded-below-tor-amplitude@cref}{{[originallemma][0][]16930}{[1][0][]source}{}{}{}}
\makeatother
\crefname{originallemma}{original lemma}{original lemmas}
\Crefname{originallemma}{Original lemma}{Original lemmas}
\makeatletter
\newlabel{lemma-diagonal-better}{{18036}{source}{}{sourcepointer-5}{}}
\newlabel{lemma-diagonal-better@cref}{{[originallemma][0][]18036}{[1][0][]source}{}{}{}}
\makeatother
\crefname{originallemma}{original lemma}{original lemmas}
\Crefname{originallemma}{Original lemma}{Original lemmas}
\makeatletter
\newlabel{lemma-internal-hom-diagonal-better}{{18799}{source}{}{sourcepointer-6}{}}
\newlabel{lemma-internal-hom-diagonal-better@cref}{{[originallemma][0][]18799}{[1][0][]source}{}{}{}}
\makeatother
\crefname{originallemma}{original lemma}{original lemmas}
\Crefname{originallemma}{Original lemma}{Original lemmas}
\makeatletter
\newlabel{lemma-perfect-ring-classical-generator}{{20288}{source}{}{sourcepointer-7}{}}
\newlabel{lemma-perfect-ring-classical-generator@cref}{{[originallemma][0][]20288}{[1][0][]source}{}{}{}}
\makeatother
\title{The diagonal tensor-internal-Hom map is an isomorphism for a perfect argument or for a pseudo-coherent argument with bounded below second object and tor amplitude condition}
\author{The Stacks Project Authors}
\date{stacks-project:more-algebra.tex:local-snapshot-20261012; GNU FDL 1.2 or later}
\begin{document}
\maketitle
\setcounter{subsection}{1}
\noindent\textit{Editorial note.} All delivered mathematical text is the author's own, copied from the pinned Stacks Project chapter file more-algebra.tex; no statement or proof step was rewritten, reordered or removed.\par\medskip
\noindent\textit{Editorial note.} The chapter file pinned-original-source.tex carries only \textbackslash{}input\{preamble\} on its first line, so no preamble text lives inside the chapter. The retained file sources/stacks-more-algebra/retained-preamble.tex is a mechanical, line-for-line copy of the project's single shared preamble as staged in the snapshot; nothing was added, removed or redefined and no package was added.\par\medskip
\noindent\textit{Editorial note.} The source URL is the source-level Stacks Project URL, as required when more than one card is produced from one source; the per-claim tag is carried by claim\_key.\par\medskip
\noindent\textit{Editorial note.} The eight references inside the counted proof span are recorded as non-load-bearing navigation pointers.\par\medskip
\noindent\textit{Allowed standard prerequisites (editorial).}\begin{itemize}
\item K-injective complexes and the derived tensor product as a total complex. used throughout the counted proof span. Supplies the explicit complexes between which the map is compared.
\item Finitely generated projective modules and the tensor-evaluation adjunction. used in the counted proof span. Is the reason each component of the displayed map is an isomorphism.
\item Boundedness and finite injective dimension as hypotheses controlling the indices. the two cases in the statement and their use in the counted proof span. Is what makes the finiteness bookkeeping work in each case.
\end{itemize}\medskip
\ifdefined\hypertarget\hypertarget{sourcepointer-1}{}\fi
\noindent\textit{Source reference.} Original definition 15874: sources/stacks-more-algebra/pinned-original-source.tex, source lines 15874--15874. Same-chapter navigation pointer; the statement it names is standard for the section and the argument that consumes it stays inside the counted proof span.\par\medskip
\ifdefined\hypertarget\hypertarget{sourcepointer-2}{}\fi
\noindent\textit{Source reference.} Original lemma 0ATH: focused-seeds/AGP/sources/stacks-snapshot-20261011/tags.txt, source lines 13615--13615. Cross-chapter navigation pointer into the Stacks Project; the standard fact it supplies is applied inside the counted proof span and the argument resting on it stays there.\par\medskip
\ifdefined\hypertarget\hypertarget{sourcepointer-3}{}\fi
\noindent\textit{Source reference.} Original lemma 18708: sources/stacks-more-algebra/pinned-original-source.tex, source lines 18708--18708. Same-chapter navigation pointer; the statement it names is standard for the section and the argument that consumes it stays inside the counted proof span.\par\medskip
\ifdefined\hypertarget\hypertarget{sourcepointer-4}{}\fi
\noindent\textit{Source reference.} Original lemma 16930: sources/stacks-more-algebra/pinned-original-source.tex, source lines 16930--16930. Same-chapter navigation pointer; the statement it names is standard for the section and the argument that consumes it stays inside the counted proof span.\par\medskip
\ifdefined\hypertarget\hypertarget{sourcepointer-5}{}\fi
\noindent\textit{Source reference.} Original lemma 18036: sources/stacks-more-algebra/pinned-original-source.tex, source lines 18036--18036. Same-chapter navigation pointer; the statement it names is standard for the section and the argument that consumes it stays inside the counted proof span.\par\medskip
\ifdefined\hypertarget\hypertarget{sourcepointer-6}{}\fi
\noindent\textit{Source reference.} Original lemma 18799: sources/stacks-more-algebra/pinned-original-source.tex, source lines 18799--18799. Same-chapter navigation pointer; the statement it names is standard for the section and the argument that consumes it stays inside the counted proof span.\par\medskip
\ifdefined\hypertarget\hypertarget{sourcepointer-7}{}\fi
\noindent\textit{Source reference.} Original lemma 20288: sources/stacks-more-algebra/pinned-original-source.tex, source lines 20288--20288. Same-chapter navigation pointer; the statement it names is standard for the section and the argument that consumes it stays inside the counted proof span.\par\medskip
\begin{lemma}
\label{lemma-internal-hom-evaluate-tensor-isomorphism}
Let $R$ be a ring. Let $K, L, M$ be objects of $D(R)$. The map
$$
K \otimes_R^\mathbf{L} R\Hom_R(M, L) \longrightarrow
R\Hom_R(M, K \otimes_R^\mathbf{L} L)
$$
of Lemma \ref{lemma-internal-hom-diagonal-better}
is an isomorphism in the following cases
\begin{enumerate}
\item $M$ perfect, or
\item $K$ is perfect, or
\item $M$ is pseudo-coherent, $L \in D^+(R)$, and $K$ has
tor amplitude in $[a, \infty]$.
\end{enumerate}
\end{lemma}

\begin{proof}
Proof in case $M$ is perfect. Note that both sides of the arrow
transform distinguished triangles in $M$ into distinguished triangles
and commute with direct sums. Hence it suffices to check
it holds when $M = R[n]$, see
Derived Categories, Remark \ref{derived-remark-check-on-generator}
and Lemma \ref{lemma-perfect-ring-classical-generator}.
In this case the result is obvious.

\medskip\noindent
Proof in case $K$ is perfect. Same argument as in the previous case.

\medskip\noindent
Proof in case (3). We may represent $K$ and $L$ by
bounded below complexes of $R$-modules $K^\bullet$ and $L^\bullet$.
We may assume that $K^\bullet$ is a K-flat complex
consisting of flat $R$-modules, see
Lemma \ref{lemma-bounded-below-tor-amplitude}.
We may represent $M$ by a bounded above complex $M^\bullet$
of finite free $R$-modules, see Definition \ref{definition-pseudo-coherent}.
Then the object on the LHS is represented by
$$
\text{Tot}(K^\bullet \otimes_R \Hom^\bullet(M^\bullet, L^\bullet))
$$
and the object on the RHS by
$$
\Hom^\bullet(M^\bullet, \text{Tot}(K^\bullet \otimes_R L^\bullet))
$$
This uses Lemma \ref{lemma-RHom-out-of-projective}.
Both complexes have in degree $n$ the module
$$
\bigoplus\nolimits_{p + q + r = n} K^p \otimes \Hom_R(M^{-r}, L^q) =
\bigoplus\nolimits_{p + q + r = n} \Hom_R(M^{-r}, K^p \otimes_R L^q)
$$
because $M^{-r}$ is finite free (as well these are finite direct sums).
The map defined in Lemma \ref{lemma-internal-hom-diagonal-better}
comes from the map of complexes defined in
Lemma \ref{lemma-diagonal-better} which uses
the canonical isomorphisms between these modules.
\end{proof}

\end{document}
